% ---------------------------------------------------------------------------------------------
% Manuscript text for the relative-vs-absolute correction argument.
%
% Every number below comes from results/statistics/relative_vs_absolute_statistics.parquet, and
% is reproduced by:
%     python -m experiments.relative_vs_absolute
%     python -m plotting.relative_vs_absolute
% The console report of the first command prints each number in the same order it appears here.
% Coverage: 19 trials, 11 subjects, 7 joints, 56,951,835 samples.
%
% Notation follows the filter: J is the parent segment, K the child, matching y^J / y^K in the
% method section and vector_sensor_stds_parent / _child in src/RelativeFilterPlus.py.
% ---------------------------------------------------------------------------------------------

\subsection{Why the correction is applied relatively}
\label{sec:relative-correction}

Two IMUs spanning a joint measure the same physical field twice, once in each sensor's body
frame, and there are two ways to turn that into an orientation measurement. A per-segment
attitude filter compares each sensor against a single global reference $\mathbf{m}_G$, giving
residuals $\hat{R}_J \mathbf{m}_J - \mathbf{m}_G$ and $\hat{R}_K \mathbf{m}_K - \mathbf{m}_G$.
The relative filter compares the two sensors against \emph{each other},
$\hat{R}_J \mathbf{m}_J - \hat{R}_K \mathbf{m}_K$, with no global reference anywhere in it.
Both residuals vanish only if the field is uniform over the two sensor sites; neither is, so
both are non-zero at the true orientation and each filter responds by rotating its estimate away
from truth until the residual it sees is nulled. We quantify how much rotation each residual
demands, and we do so from the data alone: no filter is run, and the ground-truth marker
orientations are used to place both measurements in a common frame.

For each joint, field and sample we rotate the two body-frame readings into the world frame
using the marker orientations and normalise, giving unit vectors $\mathbf{u}_J$ and
$\mathbf{u}_K$ on $S^2$, and we write $\mathbf{g}$ for the unit global reference (the subject's
median world-frame field for the magnetometer, gravity for the accelerometer; the accelerometer
is first rigid-body projected to the shared joint centre, as the filter does). Let $C_J$ be the
minimum rotation carrying $\mathbf{u}_J$ onto $\mathbf{g}$ and $C_K$ likewise for
$\mathbf{u}_K$; these are the world-frame orientation errors the per-segment route incurs, of
magnitudes $\theta_J$ and $\theta_K$. The relative residual instead demands only
$\theta_{rel} = \angle(\mathbf{u}_J,\mathbf{u}_K)$. Because the per-segment route reports a
relative orientation of $R_J^{\top} C_J^{\top} C_K R_K$, the joint angle it returns is wrong by
a conjugate of $C_{\mathrm{comp}} = C_J^{\top} C_K$, and so by the angle
$\theta_{\mathrm{comp}} = \lVert \log C_{\mathrm{comp}} \rVert$. Four exact statements follow,
each verified on every sample (Sec.~\ref{sec:relative-results}):

\begin{enumerate}
  \item \textbf{Geodesic triangle inequality.}
    $\theta_{rel} \le \theta_J + \theta_K$. Great-circle distance is a metric on $S^2$, so the
    direct route from $\mathbf{u}_J$ to $\mathbf{u}_K$ is never longer than the detour through
    $\mathbf{g}$. The relative correction never exceeds the total correction the per-segment
    route must apply across the two segments.

  \item \textbf{Minimum rotation property.}
    $\theta_{rel} \le \theta_{\mathrm{comp}}$. Since $C_K \mathbf{u}_K = \mathbf{g}$ and
    $C_J^{\top}\mathbf{g} = \mathbf{u}_J$, the composed correction $C_{\mathrm{comp}}$ carries
    $\mathbf{u}_K$ onto $\mathbf{u}_J$; any rotation doing so must turn by at least the angle
    between them. The relative residual's demand is therefore the exact lower bound on the joint
    angle error achievable by \emph{any} reference-based correction of the same pair. The
    quantifier matters: the bound holds pointwise for every $\mathbf{g}$, so it cannot be evaded
    by calibrating the global reference better.

  \item \textbf{An exact closed form.} Writing $\psi$ for the angle of
    $C_{\mathrm{comp}} C_{\min}^{\top}$, where $C_{\min}$ is the minimum rotation carrying
    $\mathbf{u}_K$ onto $\mathbf{u}_J$,
    \begin{equation}
      \cos\!\left(\tfrac{\theta_{\mathrm{comp}}}{2}\right)
        = \cos\!\left(\tfrac{\theta_{rel}}{2}\right)\cos\!\left(\tfrac{\psi}{2}\right).
      \label{eq:closed-form}
    \end{equation}
    Both rotations carry $\mathbf{u}_K$ onto $\mathbf{u}_J$, so their ratio fixes
    $\mathbf{u}_J$ and is a rotation \emph{about} it --- the per-segment route's excess is a pure
    twist about the measured field direction, which is precisely the direction a single vector
    measurement cannot observe. That twist axis is orthogonal to $C_{\min}$'s, so the half-angle
    formula for composed rotations loses its cross term and collapses to
    Eq.~\eqref{eq:closed-form}, from which statement 2 follows again with equality if and only if
    $\psi = 0$.

  \item \textbf{The excess is a spherical area.} $\psi$ is the spherical excess, equivalently the
    signed area, of the triangle $(\mathbf{u}_J, \mathbf{g}, \mathbf{u}_K)$, wrapped to
    $(-180^\circ, 180^\circ]$. The penalty for routing through a global reference is therefore
    exactly zero when that reference lies on the great circle through the two measurements, and
    grows with how far off that plane it lies.
\end{enumerate}

Nothing above is specific to magnetism, and we apply it to both fields, which fail the
uniform-field assumption for unrelated reasons. The magnetometer's $\mathbf{u}_J$ and
$\mathbf{u}_K$ depart from $\mathbf{g}$ because the field is locally distorted by ferrous
material; distortion is spatially correlated, so two sensors a segment apart see similar fields.
The accelerometer's depart from gravity because the body accelerates --- and here the mechanism
is sharper, since once both readings are projected to the shared joint centre both sensors
measure the specific force at the same physical point, so linear acceleration is exactly
common-mode and cancels from the relative residual while surviving in full in the absolute one.
We therefore carry the unprojected accelerometer through the whole analysis as a control.

\paragraph{The irreducible component.}
Statements 1--4 bound the rotation a residual demands, and a demanded rotation can cancel
between the two segments. To bound the error that \emph{cannot} be removed we use the angle
between the accelerometer and magnetometer vectors,
$\beta = \angle(\mathbf{a}, \mathbf{m})$, which is a rotation invariant: it is measured in the
sensor's own body frame, so no orientation estimate enters it and none can change it. Writing
$\beta_J$, $\beta_K$ for the two sensors and $\beta_G = \angle(\mathbf{g}, \mathbf{m}_G)$ for
the global reference pair, three further statements hold:

\begin{enumerate}\setcounter{enumi}{4}
  \item \textbf{An invariant mismatch is irreducible, and splits evenly.} For two unit
    observation pairs whose inter-vector angles differ by $\Delta\beta$, the rotation minimising
    the summed squared alignment error leaves each vector misaligned by exactly
    $\Delta\beta/2$, and no rotation does better than $\Delta\beta/2$ on the worse of the two.
    Hence $\Delta\beta/2$ is a floor on the orientation error, not a bound on one estimator.

  \item \textbf{The floors obey the same triangle inequality}, now on the real line:
    $|\beta_J - \beta_K| \le |\beta_J - \beta_G| + |\beta_G - \beta_K|$. The relative residual
    must reconcile $\beta_J$ against $\beta_K$; the absolute residual must reconcile each of them
    against $\beta_G$. Unlike statements 1--4 this survives translation into a statement about
    joint angles, because the mismatch cannot be shuffled between the two segments.

  \item \textbf{Both corrupting mechanisms are common-mode across a joint and absent from the
    reference pair.} A distorted local field tilts $\mathbf{m}$ similarly at both sensors, and
    after projection to the joint centre linear acceleration tilts $\mathbf{a}$ identically at
    both, so $\beta_J - \beta_K$ is small while $\beta_J - \beta_G$ and $\beta_K - \beta_G$ are
    not.
\end{enumerate}

This component is also the one part of the analysis independent of the marker data: $\beta$ is
computed from raw body-frame readings, so the residual sensor-to-segment alignment that inflates
$\theta_{rel}$, $\theta_J$ and $\theta_K$ cannot touch it.

\paragraph{What these quantities are not.}
They are the orientation error each residual \emph{demands} at the true orientation, not a
filter's realised error. A real filter sees both fields and a gyroscope prediction and
distributes the correction across the two segments according to its covariance; to first order at
equal trust the Kalman update takes the minimum-norm state correction consistent with the
residual, which is what makes $\theta_{rel}$ attained rather than merely a bound. Three further
limits are worth stating. First, a single vector measurement leaves rotation about its own
direction unobserved, so these are the components the measurement constrains and the gyroscope
carries the rest. Second, the marker-based alignment inflates $\theta_{rel}$ by both segments'
residual misalignment but each of $\theta_J$, $\theta_K$ by only one, so it biases every measured
gap \emph{against} the conclusion drawn here and the reported gaps are lower bounds. Third, the
relative correction is not bounded by $\min(\theta_J, \theta_K)$ and we do not claim it is: if
one sensor happens to sit on the global field its absolute correction vanishes while
$\theta_{rel}$ is the other sensor's full deviation. That case occurs on 56\% of magnetometer
samples and is reported, but it is not an implementable method, since nothing tells a
per-segment filter which of its two sensors to trust.

\subsection{Results}
\label{sec:relative-results}

All seven statements held on every one of the $5.70\times10^{7}$ samples across 19 trials, 11
subjects and 7 joints, in all 917 joint\,$\times$\,field\,$\times$\,reference cells: the worst
violation of statement~1 was $1.1\times10^{-11}$ deg and of statement~2 $9.4\times10^{-10}$ deg,
the closed form of Eq.~\eqref{eq:closed-form} was satisfied to $4.2\times10^{-14}$, and $\psi$
matched the spherical excess of statement~4 to $2.8\times10^{-6}$ deg --- floating-point noise in
every case, against a tolerance of $10^{-4}$ deg.

The two bounds differ enormously in size, and we report them separately. Across both segments,
the per-segment route costs a median of $20.5^\circ$ ($\mathrm{IQR}$ $11.9$--$40.9$) for the
magnetometer and $18.7^\circ$ ($8.5$--$40.0$) for the projected accelerometer, against
$7.9^\circ$ ($4.4$--$14.8$) and $5.6^\circ$ ($2.5$--$11.9$) for the relative correction: gaps of
$12.6^\circ$ and $13.1^\circ$ at the median, with the relative correction smaller in 131/131
trial\,$\times$\,joint cells ($p = 1.5\times10^{-23}$, Wilcoxon). In the \emph{joint angle},
however, the same comparison gives $7.87^\circ$ versus $7.85^\circ$ and $5.65^\circ$ versus
$5.61^\circ$ --- gaps of $0.02^\circ$ and $0.04^\circ$. This is a negative result and we report
it as one: correcting both sensors to a shared global reference produces two orientation errors
that are largely common-mode and cancel in the joint angle between them.
Equation~\eqref{eq:closed-form} explains it quantitatively, since for small angles it reduces to
$\theta_{\mathrm{comp}} \approx \sqrt{\theta_{rel}^2 + \psi^2}$, so the out-of-plane twist adds
in \emph{quadrature}; with a median $|\psi|$ of $0.23^\circ$ (magnetometer) and $0.20^\circ$
(accelerometer) against a $\theta_{rel}$ of several degrees, it costs almost nothing. The
measured excess tracks $|\psi|$ closely (Spearman $\rho = 0.96$ and $0.98$ over the per-cell
medians) and follows Eq.~\eqref{eq:closed-form} across six orders of magnitude in the excess
with no fitted parameter (Fig.~\ref{fig:relative-vs-absolute}C). None of this depends on how the global
reference is chosen: substituting the torso-only field the per-segment baseline actually uses,
or a per-joint least-squares-optimal constant direction, changes the gap by less than
$0.02^\circ$.

The joint-angle advantage instead appears in the irreducible component. The acc--mag angle was
$149.6^\circ$ at the parent sensor and $147.8^\circ$ at the child against $152.8^\circ$ for the
global reference pair, giving a relative mismatch of $6.9^\circ$ ($3.0$--$14.2$) against absolute
mismatches of $8.1^\circ$ and $10.3^\circ$. By statement~5 the corresponding floors are
$3.5^\circ$ ($1.5$--$7.1$) for the relative route and $6.5^\circ$ ($3.4$--$11.7$) for the
per-segment route: a factor of $1.9$, or $3.1^\circ$ of orientation error that no choice of
orientation can remove. The relative floor was smaller in all 44 subject\,$\times$\,joint blocks
(rank-biserial $-1.00$, Holm-corrected Wilcoxon $p < 10^{-4}$; blocked by subject and joint type,
averaged over side and activity), and the ratio held at every joint, from $1.9\times$ at the
lumbar joint to $1.6$--$2.2\times$ at the ankles and knees. Removing the joint-centre projection
raises the relative floor from $3.5^\circ$ to $4.5^\circ$ ($p < 10^{-4}$, rank-biserial $-0.97$)
while moving the absolute floor only from $6.5^\circ$ to $6.9^\circ$, confirming that pairing the
sensors \emph{at a common physical point} --- and not merely pairing them --- is what makes the
relative residual consistent.

% ---------------------------------------------------------------------------------------------
% Figure caption
% ---------------------------------------------------------------------------------------------

\begin{figure}[t]
  \centering
  \includegraphics[width=\textwidth]{relative_vs_absolute}
  \caption{\textbf{A relative correction is bounded below a global-reference correction, and its
    irreducible part is half the size.}
    \textbf{(A)} The geometry. The two world-frame measurement directions $\mathbf{u}_J$,
    $\mathbf{u}_K$ and the global reference $\mathbf{m}_G$ form a spherical triangle whose area
    is $\psi$; Eq.~\eqref{eq:closed-form} makes the per-segment route's excess an exact function
    of $\psi$, which vanishes when $\mathbf{m}_G$ lies on the great circle through the two
    measurements (dashed). Angles are exaggerated for legibility.
    \textbf{(B)} Both bounds against the identity line. The joint-angle bound
    ($\theta_{\mathrm{comp}}$) is tight --- the two per-segment errors are common-mode and cancel
    --- while the both-segments bound ($\theta_J + \theta_K$) is not.
    \textbf{(C)} The measured excess against the out-of-plane twist, in bands of $\theta_{rel}$,
    with Eq.~\eqref{eq:closed-form} overlaid (dashed, no fitted parameters).
    \textbf{(D,\,E)} Correction angles by joint for the magnetometer and for the accelerometer
    after projection to the joint centre: the same geometry driven by field distortion and by
    rigid-body kinematics respectively.
    \textbf{(F)} Pooled distributions; $\theta_{rel}$ and $\theta_{\mathrm{comp}}$ coincide.
    \textbf{(G)} The irreducible floors by joint, half the acc--mag angle mismatch, with the
    unprojected control (open triangles).
    \textbf{(H)} The same per subject\,$\times$\,joint cell against the identity line.
    \textbf{(I)} One trial at full rate through its busiest window: the two per-segment angles
    rise and fall together with each stride. Boxes show the IQR with p5--p95 whiskers.
    Panels B--H pool $4.9\times10^{6}$ stored samples from 19 trials, 11 subjects and 7 joints;
    ground-truth orientations are from markers and no filter is run.}
  \label{fig:relative-vs-absolute}
\end{figure}
